% GENERATED FILE. Edit canonical lessons, then run npm run generate:books.
% canonical-source-hash: fnv1a64:2c54a755a60c30bc
% canonical-lessons: KA-C38-hrudaya, KA-S132-sign-visarga

\chapter{The Heart}
\label{ch:ka-heart}
\hlchaptermodality{38}

\begin{chapteropening}
\textbf{By the end of this chapter:} \emph{I can name my heart in Kannada and say why it is a genuine Indo-European cousin of the English word, not just a resemblance.}

Say \kn{ನನ್ನ} \kn{ಹೃದಯ}, read the vocalic-r sign \kn{ೃ} against Chapter 14's independent \kn{ಋ}, and trace hṛdaya to the same Proto-Indo-European root as English heart, Latin cor and Greek kardía.
\end{chapteropening}

\section[hṛdaya]{\kn{ಹೃದಯ} (hṛdaya) — \textquotedblleft{}heart,\textquotedblright{} a real cousin of the English word}

\label{lesson:KA-C38-hrudaya}



\begin{quote}
Six body words so far, all Dravidian. This one is Sanskrit — and it hides the closest kinship to English that this whole book has found.
\end{quote}

\subsection*{You'll want to know: \kn{ಹೃದಯ}}
\begin{quote}
\textbf{\kn{ಹೃದಯ}} — \emph{hṛdaya} — \textbf{heart}
\end{quote}

\begin{quote}
\textbf{\kn{ನನ್ನ} \kn{ಹೃದಯ}.} — \emph{nanna hṛdaya.} — \textquotedblleft{}My heart.\textquotedblright{}
\end{quote}

\subsection*{The letters in this word}
\emph{(Skip this if you already read Kannada.)} \textbf{\kn{ಹೃ}} is \textbf{\kn{ಹ}} (\emph{ha}) carrying a sign you have not met: \textbf{\kn{ೃ}}, the \textbf{vocalic-r} vowel sign, said \emph{ṛ} — a short, tapped \textquotedblleft{}ri.\textquotedblright{} Chapter 14 taught this same sound as a full \textbf{independent} letter, \textbf{\kn{ಋ}}, at the front of \textbf{\kn{ಋತು}} (\emph{ṛtu}, \textquotedblleft{}season\textquotedblright{}), because a word can only open with the independent form. Here the sound falls in the middle of a word, so it rides on a consonant as the small sign \textbf{\kn{ೃ}} instead: \textbf{\kn{ಹ} + \kn{ೃ} = \kn{ಹೃ}}, \emph{hṛ}. Same sound, two different shapes, depending only on where in the word it lands.

\begin{cousinweb}
\textbf{\kn{ಹೃದಯ}} is Sanskrit \textbf{\dv{हृदय}} (\emph{hṛdaya}), from Proto-Indo-European \textbf{*\'{k}ērd-}. That root also gave English \textbf{heart}, Latin \textbf{cor} (behind English \emph{courage} and \emph{cordial}), and Greek \textbf{kardía} (behind \emph{cardiac}). Every earlier cross-family match in this book — \emph{iṣṭa} and \emph{ask}, \emph{sahāya} and \emph{same}, Spanish's \emph{me gusta} — was two languages landing on the same idea by accident, unrelated in ancestry. This one is different: Sanskrit and English are both \textbf{Indo-European}, so \emph{hṛdaya} and \emph{heart} are not a coincidence. They are real cousins, descended from the very same word.
\end{cousinweb}

\subsection*{Your turn}
\begin{itemize}
\raggedright
  \item \emph{Say it:} \textquotedblleft{}nanna hṛdaya\textquotedblright{} — my heart
  \item \emph{Say it:} the vocalic-r, word-middle and word-initial — \textquotedblleft{}hṛ … ṛtu\textquotedblright{}
  \item \emph{Say it:} the real cousins — \textquotedblleft{}hṛdaya … heart … cor … kardía\textquotedblright{}
  \item \emph{Say it:} seven body words now — \textquotedblleft{}tale, kai, kaṇṇu, kivi, mūgu, bāyi, hṛdaya\textquotedblright{}
  \item \emph{Recall:} say \emph{kuṭumba}
\end{itemize}

\begin{tcolorbox}[breakable,colback=teal!4,colframe=teal!35!black,title={Before you move on}]
Say \textquotedblleft{}my heart.\textquotedblright{} (\emph{Nanna hṛdaya}.) Where else have you seen the vocalic-\emph{r} sound, and in what shape? (\textbf{Chapter 14's} \emph{ṛtu} — as the \textbf{independent} letter \kn{ಋ}, word-initially.) Is \emph{hṛdaya}'s match with English \emph{heart} an accident, like \emph{iṣṭa} and \emph{ask}? (\textbf{No} — Sanskrit and English are both \textbf{Indo-European}; this is real shared ancestry.) Name two other English words from the same root. (\textbf{Cardiac}, from Greek; \textbf{courage} or \textbf{cordial}, from Latin.)
\end{tcolorbox}
\section[visarga]{\kn{◌ಃ} — one character, two dots, a breath at the end}

\label{lesson:KA-S132-sign-visarga}



\begin{quote}
Before the new one: \kn{◌ೃ} — where does it sit, and which standing letter is it the other half of?

One character this time. It has a partner you met long ago, and the two are easiest to learn side by side.
\end{quote}

\subsection*{Script you'll notice: \kn{◌ಃ}}
\textbf{\kn{◌ಃ}} — two dots, one above the other, sitting after a syllable.

Kannada has \textbf{two marks that ride at the end of a syllable} rather than standing as letters. You already read one:

\begin{itemize}
\raggedright
  \item \textbf{\kn{◌ಂ}} — one small circle. Turns the end of the syllable \textbf{nasal}, the sound in the middle of \textbf{\kn{ಸಂತೋಷ}} \emph{santōṣa}.
  \item \textbf{\kn{◌ಃ}} — two dots. Adds a \textbf{breath}, an \emph{h} let out after the vowel.
\end{itemize}

Say \emph{santōṣa} and feel the air go up into your nose on the \textbf{\kn{◌ಂ}}. Now say \textbf{ah} and let it fall away as a soft \emph{ha}. That falling breath is what the two dots are for.

Hold the two marks as a pair: neither is a letter you can sound out alone, and both attach to what came before them.

\textbf{Where you will meet it.} Rarely, and almost always ending a word Kannada borrowed from Sanskrit. Knowing the two dots on sight is enough.

\subsection*{Writing: \kn{◌ಃ} — copy what you see}
Put your pen on \kn{◌ಃ} and follow its line. Copy the shape you can see — slowly, and larger than it is printed.

\begin{quote}
This book does not yet tell you \textbf{where to start the character or which way to travel}. That is a real thing, taught with real variation from school to school, and it is not written down here until it can be written down with a source. Copying what is in front of you needs no such source, and it is how the shape gets into your hand in the meantime.
\end{quote}

\subsection*{Your turn}
\begin{itemize}
\raggedright
  \item \emph{Look:} at these two marks and say which one turns the syllable nasal
\end{itemize}

\begin{quote}
\kn{◌ಂ}  ·  \kn{◌ಃ}
\end{quote}

\begin{itemize}
\raggedright
  \item \emph{Trace it:} \kn{◌ಃ} three times, letting a soft breath out as you finish each one
  \item \emph{Look:} at \kn{ಸಂತೋಷ} and find the other mark inside it
  \item \emph{Say it:} the nasal one, then the breathed one, and name what each adds
\end{itemize}

\begin{tcolorbox}[breakable,colback=teal!4,colframe=teal!35!black,title={Before you move on}]
Which character is this — \kn{◌ಃ}, and what does it add to the syllable before it? (\textbf{A breath, an \emph{h} after the vowel}.) Which mark is its partner, and what does that one add? (\textbf{\kn{◌ಂ}}, a \textbf{nasal}.)
\end{tcolorbox}
